\section{Display dan Layout Dasar}

Properti \texttt{display} mengontrol bagaimana elemen ditampilkan dan berpartisipasi dalam layout \cite{mdn-learn-css}:
\begin{itemize}
  \item \texttt{block}: elemen mengambil lebar penuh dan dimulai baris baru (contoh: \texttt{div}, \texttt{p})
  \item \texttt{inline}: elemen mengalir dalam baris yang sama, lebar sesuai konten (contoh: \texttt{span}, \texttt{a})
  \item \texttt{inline-block}: seperti inline tetapi dapat diberi width/height dan margin/padding vertikal
\end{itemize}

\textbf{Positioning}: \texttt{position} dengan nilai \texttt{static} (default), \texttt{relative}, \texttt{absolute}, \texttt{fixed}, atau \texttt{sticky} mengontrol letak elemen. \texttt{relative} menggeser elemen relatif terhadap posisi awalnya; \texttt{absolute} memosisikan relatif terhadap ancestor yang ber-\texttt{position} non-static \cite{w3schools-css}.

Warna dan teks: properti seperti \texttt{color}, \texttt{background-color}, \texttt{font-family}, \texttt{font-size}, \texttt{font-weight}, \texttt{text-align}, \texttt{line-height} digunakan untuk mengatur tampilan teks dan latar \cite{mdn-learn-css}. Satuan yang umum: \texttt{px}, \texttt{em}, \texttt{rem}, \texttt{\%}.
